\chapter{赛道与策略：如何定义你的「第一只基金」}

\section{聚焦优于泛化}

首只基金若宣称「全行业、全阶段、全球配置」，在 allocator 视角中往往等同于\textbf{无清晰边际优势}。更可行的做法是在\textbf{行业 × 地域 × 阶段 × 单笔规模}四维中至少锁定两维，形成可辩护的「为什么是你」的故事。聚焦不等于永久狭窄：可在 LPA 中保留\textbf{合理扩展投资范围}的修正案机制，但首轮沟通应强调核心圈。

\section{可重复的筛选标准}

建立书面化的\textbf{立项标准}：最低收入/利润门槛、市场集中度上限、禁止行业（如强监管敏感领域）、ESG 红线、以及技术与团队评分卡。标准应可回溯：每个过会/否决决策对应哪一条标准被触发。这样既能训练初级投资人员，也能向 LP 证明流程非拍脑袋。量化背景团队可将部分标准\textbf{指标化}，但仍需与投资委员会的人文判断结合。

\section{基金规模与机会集的匹配}

目标募资额应与\textbf{可完成投资节奏}及\textbf{团队覆盖能力}一致：过小则管理费无法覆盖固定成本；过大则被迫降低标准或拉长投资期，损害 IRR 与声誉。冷启动阶段可采用较低硬上限 + 后续基金接力，而不是单支基金「一步到位」。

\section{品牌与执行的乘数}

赛道定义会通过每一次公开露面、案例研究与 LP 推荐被\textbf{重复强化}。频繁切换叙事会稀释品牌，使介绍链无法累积。建议每年复盘一次策略声明，仅在证据充分时调整表述，并同步修订 PPM 与网站（若有）。
